%% =========================================================================
%% Section 7: Methodological Sensitivity and Robustness Analyses
%% =========================================================================

\section{Methodological Sensitivity and Robustness Analyses}

To establish that our empirical findings are not artifacts of specific analytical choices, variable cutoffs, or sample definitions, we subjected all models to an extensive suite of sensitivity and robustness tests.

\subsection{Response Modeling and Inverse-Probability Weighting (IPW)}
As documented in the methodology section, 7.8\% of eligible non-proxy respondents had incomplete direct cognitive assessments, with missingness concentrated among adults aged 75+ (19.9\%) and unschooled respondents (13.3\%). If individuals with missing cognitive scores had systematic, unobserved cognitive deficits correlated with early childhood poverty, complete-case analyses could potentially suffer from selection bias.

To address this, we estimated a multivariable survey-weighted logistic response model predicting complete cognitive assessment among all direct interviewees ($N = 66{,}606$) as a function of childhood exposures, demographic characteristics, state fixed effects, adult education, consumption expenditure, and physical health conditions (Table~\ref{tab:ipw_response}).

\begin{table}[!htbp]
\centering
\small
\caption{\textbf{Multivariable Logistic Model Predicting Complete Cognitive Assessment ($\mathbf{N=66{,}606}$).}}
\label{tab:ipw_response}
\resizebox{0.88\textwidth}{!}{%
\begin{tabular}{lccc}
\toprule
\textbf{Predictor Variable} & \textbf{Comparison Level (vs. Reference)} & \textbf{Odds Ratio (OR)} & \textbf{\textit{p}-value} \\
\midrule
Gender & Female (vs. Male) & 0.903 & $0.003$ \\
Sector & Urban (vs. Rural) & 1.188 & $<0.001$ \\
Social Group (Caste) & Scheduled Tribe (vs. General/None) & 0.718 & $<0.001$ \\
                     & Scheduled Caste (vs. General/None) & 1.021 & $0.693$ \\
                     & Other Backward Classes (vs. General) & 1.017 & $0.711$ \\
\addlinespace
Educational Attainment & Less than primary (vs. None) & 1.624 & $<0.001$ \\
                       & Primary completed (vs. None) & 2.025 & $<0.001$ \\
                       & Middle completed (vs. None)  & 2.289 & $<0.001$ \\
                       & Secondary completed (vs. None) & 3.277 & $<0.001$ \\
                       & Higher secondary/Diploma (vs. None) & 3.574 & $<0.001$ \\
                       & Graduate and above (vs. None) & 4.168 & $<0.001$ \\
\addlinespace
Household MPCE Quintile & Quintile 2 (vs. Quintile 1, Poorest) & 1.063 & $0.165$ \\
                        & Quintile 3 (vs. Quintile 1)          & 1.135 & $0.005$ \\
                        & Quintile 4 (vs. Quintile 1)          & 1.347 & $<0.001$ \\
                        & Quintile 5 (vs. Quintile 1, Richest) & 1.381 & $<0.001$ \\
\addlinespace
General Physical Health & Poor Self-Rated Health (vs. Excellent) & 0.567 & $<0.001$ \\
Childhood SES           & Poor Childhood SES (vs. Well-off)    & 1.148 & $0.026$ \\
\bottomrule
\end{tabular}%
}
\end{table}

The response model demonstrates that missingness is \textbf{strongly informative}: higher education increases test completion odds up to 4-fold ($\text{OR} = 4.17, p < 0.001$), while Scheduled Tribe respondents ($\text{OR} = 0.72, p < 0.001$) and individuals in poor health ($\text{OR} = 0.57, p < 0.001$) are substantially less likely to complete the battery. 

To correct for this non-random response pattern, inverse-probability weights ($w_i^{\text{IPW}}$) were calculated and trimmed at the 99th percentile to prevent extreme weights, and applied to re-estimate primary Model B (Table~\ref{tab:ipw}).

\begin{table}[!htbp]
\centering
\small
\caption{\textbf{Comparison of Complete-Case vs. Inverse-Probability-Weighted (IPW) Model B Cognitive Estimates ($\mathbf{N=57{,}245}$).}}
\label{tab:ipw}
\resizebox{0.85\textwidth}{!}{%
\begin{tabular}{lcc}
\toprule
\textbf{Childhood Exposure Parameter} & \textbf{Complete-Case APR (95\% CI)} & \textbf{IPW-Adjusted APR (95\% CI)} \\
\midrule
Childhood SES: Average & 0.797 (0.703 -- 0.905) & 0.799 (0.708 -- 0.903) \\
Childhood SES: Poor & \textbf{0.928 (0.818 -- 1.052)} & \textbf{0.928 (0.821 -- 1.049)} \\
\midrule
Childhood health: Fair/Poor/Very poor & 0.949 (0.820 -- 1.099) & 0.942 (0.811 -- 1.094) \\
Childhood health: Poor/Very poor (orig) & 1.070 (0.813 -- 1.408) & 1.051 (0.801 -- 1.378) \\
Childhood bedridden $\ge$1 month & 1.005 (0.860 -- 1.176) & 1.003 (0.861 -- 1.169) \\
\bottomrule
\end{tabular}%
}
\end{table}

As shown in Table~\ref{tab:ipw}, the complete-case and IPW-adjusted effect estimates are \textbf{virtually identical down to the third decimal place} (Poor childhood SES $\text{APR} = 0.928$ in both specifications). This demonstrates that under the standard Missing at Random (MAR) assumption, differential test completion does not distort the substantive conclusion that early childhood SES associations with memory-orientation impairment attenuate to the null.

\subsection{The Cognitive Battery Conundrum: Sensitivity to Test Construction}
A pivotal methodological finding of this study is that conclusions regarding cognitive attenuation in low- and middle-income countries depend critically on \textbf{how the cognitive battery is constructed}. Table~\ref{tab:s4a} compares childhood SES associations across four alternative definitions of the cognitive outcome.

\begin{table}[!htbp]
\centering
\small
\caption{\textbf{Sensitivity to Cognitive Outcome Definition: Prevalence Ratios (95\% CI) for Poor Childhood SES (vs. Well-off).}}
\label{tab:s4a}
\resizebox{\textwidth}{!}{%
\begin{tabular}{lcc}
\toprule
\textbf{Cognitive Outcome Definition} & \textbf{Model A PR (95\% CI)} & \textbf{Model B PR (95\% CI)} \\
\midrule
\textbf{Primary Battery:} Lowest weighted decile, mem+orient (0--28) & 1.356 (1.156--1.590)$^{***}$ & \textbf{0.928 (0.818--1.052)} \\
\textbf{Skip-Aware Full Battery:} Lowest decile (\% of max administered) & 1.970 (1.604--2.420)$^{***}$ & \textbf{1.294 (1.068--1.569)$^{**}$} \\
\textbf{Schooling-Stratum Decile:} Mem+orient, stratum-specific decile & 0.903 (0.781--1.045) & \textbf{1.001 (0.858--1.167)} \\
\textbf{Original Composite:} Literacy-unadjusted composite (skips scored 0) & 1.653 (1.343--2.035)$^{***}$ & \textbf{1.044 (0.891--1.222)} \\
\bottomrule
\multicolumn{3}{l}{\footnotesize $^{***}\ p < 0.001$, $^{**}\ p < 0.01$. Model A: age, sex, residence, caste, state FE. Model B: + adult education and MPCE.} \\
\end{tabular}%
}
\end{table}

The comparisons in Table~\ref{tab:s4a} yield vital insights into measurement confounding:
\begin{enumerate}
    \item \textbf{Residual Confounding in Full Battery:} When complex arithmetic (serial 7s, backward counting) and graphic drawing tasks are included (row 2), poor childhood SES remains significantly associated with impairment ($\text{APR} = 1.294$, 95\% CI: 1.068--1.569, $p = 0.009$) even after adult education adjustment. This occurs because pencil-and-paper and mental calculation tasks capture procedural literacy and test familiarity that standard categorical education controls cannot fully absorb.
    \item \textbf{Complete Equalization in Schooling-Stratum Deciles:} When cognitive impairment is defined relative only to peers with identical schooling (row 3), the childhood SES penalty disappears entirely even in Model A ($\text{PR} = 0.903, p = 0.171$) and is precisely null in Model B ($\text{PR} = 1.001, p = 0.991$). This confirms that childhood socioeconomic disparities in cognition operate almost exclusively through differential access to schooling.
    \item \textbf{Primary Battery Appropriateness:} Our primary memory-and-orientation metric (row 1) avoids both the artifactual confounding of arithmetic tests and the insensitivity of within-stratum definitions, offering a fair and culturally robust measure of neuro-cognitive health.
\end{enumerate}

\subsection{Stratification by Adult Schooling Status (Schooling-Specific Models)}
To isolate whether childhood poverty leaves cognitive traces independent of formal educational attainment, we stratified the common analytic sample into respondents who never attended school versus those with any formal schooling.

\begin{enumerate}
    \item \textbf{Adults with No Formal Education ($n = 25{,}684$):} Among unschooled elders, poor childhood SES was not associated with binary cognitive impairment ($\text{APR} = 0.909$, 95\% CI: 0.788--1.048, $p = 0.188$). However, on the continuous memory-orientation scale, poor childhood SES was associated with a \textbf{statistically significant memory deficit} of \textbf{$-0.386$ points} (95\% CI: $-0.725$ to $-0.047$, $p = 0.026$). This reveals that among individuals who never had the protective benefit of formal schooling, early poverty leaves a subtle, persistent deficit in episodic recall capacity.
    \item \textbf{Adults with Formal Schooling ($n = 31{,}561$):} Among adults who attended school, poor childhood SES showed no association with binary cognitive impairment ($\text{APR} = 0.992$, 95\% CI: 0.761--1.292, $p = 0.950$) and no significant deficit on the continuous scale ($\beta = -0.183$, 95\% CI: $-0.469$ to $+0.102$, $p = 0.208$). Acquiring formal education effectively insulated individuals from childhood socioeconomic deficits across all cognitive metrics.
\end{enumerate}

\subsection{Modified Poisson vs. Logistic Regression: Preventing Risk Exaggeration}
To illustrate why modified Poisson regression was methodologically required, we compared our primary Poisson Prevalence Ratios (PR) with Adjusted Odds Ratios (AOR) from conventional logistic regression:
\begin{itemize}
    \item For depressive symptoms (weighted prevalence 28.2\%), poor childhood health produced an \textbf{Adjusted Odds Ratio of 1.33} (95\% CI: 1.12--1.58).
    \item In contrast, the modified Poisson model produced an \textbf{Adjusted Prevalence Ratio of 1.20} (95\% CI: 1.07--1.34).
\end{itemize}
Because more than a quarter of the population has depressive symptoms, the Odds Ratio overstates the relative risk elevation by 65\% ($+33\%$ vs. $+20\%$). Modified Poisson regression provides policymakers with the true, uninflated relative increase in disease burden.

\subsection{Additional Robustness Checks}
All substantive conclusions remained invariant across several additional checks:
\begin{itemize}
    \item \textbf{Continuous Depression Scale:} Using continuous CES-D scores (0--10) in survey linear regressions, poor childhood health was associated with a $+0.365$ point elevation ($p < 0.001$), poor childhood SES with a $+0.324$ point elevation ($p < 0.001$), and bedridden illness with a $+0.210$ point elevation ($p = 0.004$).
    \item \textbf{Exclusion of Centenarians:} Re-running models after excluding 58 centenarians (aged $\ge 100$) produced identical estimates for depression ($\text{APR} = 1.207$, 95\% CI: 1.031--1.414, $p = 0.019, N = 62{,}011$) and cognition ($\text{APR} = 0.928$, 95\% CI: 0.818--1.052, $p = 0.242, N = 57{,}226$).
    \item \textbf{Varied Childhood SES:} Retaining the 394 respondents who reported that their childhood finances ``varied'' as a distinct category yielded $\text{APR} = 1.180$ (95\% CI: 0.889--1.566, $p = 0.253$) for depression and $\text{APR} = 1.392$ (95\% CI: 0.933--2.076, $p = 0.105$) for cognition, with primary coefficients remaining unchanged.
    \item \textbf{Original Childhood Health Coding:} Applying the narrow published dichotomy (Poor/Very poor [1.6\%] vs. the rest [98.4\%]) yielded a stronger association with depression ($\text{APR} = 1.42$, 95\% CI: 1.21--1.66, $p < 0.001$) and remained null for cognition ($\text{APR} = 1.07, p = 0.629$).
\end{itemize}
